\chapter{Special Populations}\label{ch:special}

\section{Pediatric First Aid (Infants and Children)}

Children are not simply small adults -- their physiology requires modified approaches to assessment and treatment. Common differences include higher metabolic rates, proportionally larger heads relative to bodies, different normal vital signs, and varying susceptibility to injury patterns.

\subsection{Infant Assessment (< 1 Year)}

\begin{itemize}[leftmargin=*]
  \item \textbf{Responsiveness}: Gently tap the sole of the foot or rub the back/shoulders (do not shake an infant)
  \item \textbf{Breathing}: Look for chest movement; listen for breath sounds near mouth and nose simultaneously (infants breathe primarily through nose)
  \item \textbf{Pulse}: Check brachial artery (inner upper arm) — place two fingers on the inside of the elbow crease and move slightly toward armpit to feel pulse
  \item \textbf{Normal Vital Ranges}: Heart rate 100-160 bpm (infants slow down with age), respiratory rate 30-60 breaths/min (infants), temperature >100.4 degrees Fahrenheit (38 degrees Celsius) in infants under 3 months = medical emergency
\end{itemize}

\noindent\textbf{\textcolor{red}{WARNING}}: Any rectal temperature >=100.4 degrees Fahrenheit (38 degrees Celsius) in an infant under 3 months old constitutes a medical emergency requiring immediate evaluation. Infant immune systems are immature and serious infections can progress rapidly.

\subsection{Pediatric CPR Considerations}

\begin{itemize}[leftmargin=*]
  \item Use softer compressions for infants and smaller children (approximately one-third the diameter of chest or about 1.5 inches for infant, 2 inches for child)
  \item For rescue breaths, use mouth-to-mouth-and-nose technique for infants (seal over both mouth and nose)
  \item Be cautious with head positioning — infants' heads are large relative to body; overextension of the neck during airway management can obstruct the airway. Neutral or slightly sniffing position is optimal for infants
  \item Pediatric cardiac arrest most often results from respiratory failure followed by bradycardia and arrest (unlike adults where cardiac causes predominate). Emphasize effective ventilation along with compressions
\end{itemize}

\subsection{Choking in Infants and Children}

See Chapter~\ref{ch:cpr} for detailed choking management techniques specific to each age group. Key differences:

\begin{itemize}[leftmargin=*]
  \item Never perform abdominal thrusts on an infant under 1 year
  \item Use alternating back blows and chest thrusts for infants
  \item Be gentle but firm — infants' internal organs are more vulnerable to injury
\end{itemize}

\subsection{Febrile Seizures in Children}

Brief febrile seizures occur in some young children (typically 6 months to 5 years) when fever rises quickly. They appear alarming but usually cause no long-term harm.

\begin{itemize}[leftmargin=*]
  \item Lay child on side on safe surface (prevents aspiration)
  \item Do NOT restrain or put anything in the mouth
  \item Time the seizure
  \item After seizure ends, seek medical attention for fever evaluation
  \item Emergency services if first-time seizure, lasts longer than 5 minutes, second seizure occurs before recovery, or breathing difficulty persists after seizure ends
\end{itemize}

\section{Geriatric First Aid (Older Adults)}

Aging brings physiological changes that affect both presentation of emergencies and response to interventions. Common considerations include polypharmacy, comorbidities, altered pain perception, and slower healing.

\subsection{Common Geriatric Vulnerabilities}

\begin{itemize}[leftmargin=*]
  \item Higher risk of falls and fractures due to balance issues, muscle weakness, and bone density loss
  \item Medication interactions can alter symptom presentation (e.g., beta-blockers mask tachycardia)
  \item Decreased thirst sensation leads to dehydration risks even at room temperature
  \item Fever may be blunted or absent despite serious infection
  \item Cognitive impairment (dementia, delirium) can delay recognition of symptoms or ability to communicate
  \item Skin fragility increases risk of bruising, tearing, and pressure injuries
\end{itemize}

\subsection{Stroke in Older Adults}

Symptoms may be subtle or mistaken for other age-related conditions. Always evaluate new-onset neurological changes promptly. Watch for:

\begin{itemize}[leftmargin=*]
  \item Unexplained confusion or sudden change in mental status
  \item Sudden vision changes in one or both eyes
  \item Loss of balance or coordination (difficulty walking)
  \item Sudden severe headache
  \item Difficulty speaking or understanding speech
\end{itemize}

\noindent\textbf{\textcolor{red}{MEDICATION ALERT}}: Many older adults take anticoagulants (warfarin, aspirin, clopidogrel, newer oral anticoagulants). Even minor head trauma in these patients carries increased risk of intracranial hemorrhage. Any head injury in a patient on blood thinners warrants immediate medical evaluation.

\subsection{Falls Management}

\begin{itemize}[leftmargin=*]
  \item Assess responsiveness and level of consciousness first
  \item Ask about pain location before attempting movement
  \item If injured in way that suggests possible fracture or spinal involvement, do NOT attempt to move them unless in immediate danger
  \item Check pulse and breathing; provide CPR if needed
  \item Reassure the person while waiting for emergency services — anxiety can worsen blood pressure fluctuations
  \item Once moved, check for visible injuries, especially head and hip areas (hip fractures are common and may be minimally painful initially)
  \item Monitor closely as they may deteriorate rapidly due to underlying medical conditions
\end{itemize}

\section{First Aid During Pregnancy}

Pregnancy alters anatomy and physiology, affecting how emergencies manifest and how treatments should be modified.

\subsection{Positioning Considerations}

After 20 weeks gestation, the enlarged uterus can compress the inferior vena cava when lying flat on the back (supine hypotensive syndrome), reducing blood return to the heart and causing low blood pressure.

\begin{itemize}[leftmargin=*]
  \item Position pregnant women on their left side if possible during procedures or transport \cite{ilcor_2025}
  \item If manual intervention is needed (CPR, intubation), place a rolled towel or blanket under the right hip/left side to tilt uterus off major vessels
  \item For CPR in later pregnancy, hand placement may need to be slightly higher due to breast enlargement and reduced chest space
\end{itemize}

\subsection{Labor and Delivery Emergencies}

If delivery seems imminent (woman cannot reach a hospital):

\begin{enumerate}[leftmargin=*]
  \item Place the woman on her back with legs elevated if possible
  \item Prepare clean delivery area with clean towels/cloths
  \item Support the baby's head as it crowns (emerges) — DO NOT pull on the baby
  \item Let the natural contractions deliver the baby slowly
  \item When fully delivered, clear the baby's mouth/nose gently with suction bulb or finger (if available)
  \item Dry and stimulate the baby (rub back/feet)
  \item Keep mother and baby skin-to-skin for warmth
  \item Cut the umbilical cord only if necessary (use clean scissors/tied with sterile string if absolutely required); leave baby attached until professional help arrives
  \item Monitor for excessive bleeding from the mother — apply gentle pressure to uterus below navel if soft/bloated
\end{enumerate}

\noindent\textbf{\textcolor{red}{PLACENTA PREVIA}}: If bright red vaginal bleeding occurs before labor without contractions, this may indicate placenta previa covering the cervix. Do NOT perform vaginal exam. Keep woman on side, seek emergency transport immediately.

\section{Persons with Disabilities}

Providing appropriate first aid requires understanding various disabilities and adapting communication and care accordingly.

\subsection{Visual Impairments}

\begin{itemize}[leftmargin=*]
  \item Announce yourself clearly when approaching
  \item Ask before providing physical assistance — some prefer to be guided rather than led
  \item Describe actions you are taking as you take them ("I'm going to call emergency services now," "I'm checking your pulse")
  \item Offer verbal descriptions of environment changes
  \item Respect guide animals — never pet, feed, or distract a service animal while working
\end{itemize}

\subsection{Hearing Impairments}

\begin{itemize}[leftmargin=*]
  \item Get their attention first (tap shoulder, wave within visual range)
  \item Speak facing them so they can read lips (enunciate clearly but naturally, don't exaggerate)
  \item Use gestures pointing to items/actions
  \item Write messages if necessary
  \item Identify yourself in writing if sign language interpreter is not present
  \item Confirm understanding by asking them to repeat instructions or actions
\end{itemize}

\subsection{Mobility Impairments}

\begin{itemize}[leftmargin=*]
  \item Recognize that limited mobility does not necessarily limit consciousness or cognitive function
  \item Wheelchairs are part of the person's body — keep the person seated whenever possible during assessment
  \item Remove seat belts carefully if needed during extrication; support the wheelchair during transfers
  \item If moving is necessary, use proper lifting techniques involving more people if heavy equipment involved
  \item Prosthetic limbs should generally stay attached unless causing obstruction to life-saving care
\end{itemize}

\subsection{Cognitive and Developmental Disabilities}

\begin{itemize}[leftmargin=*]
  \item People with autism, Down syndrome, or developmental delays may not understand emergency procedures or communicate distress typically
  \item Watch for non-verbal cues of pain or discomfort (grimacing, withdrawal, vocalizations)
  \item Use simple, direct language; demonstrate what you're doing
  \item Allow extra time for processing information and responding
  \item Know routine medications/supplies these individuals rely on (communication boards, sensory tools, special diets)
  \item Contact caregivers/family if present to obtain medical history information
\end{itemize}

\subsection{Medical Alert Identification}

Many individuals with significant medical conditions wear medical alert bracelets, necklaces, or carry cards. These typically contain:

\begin{itemize}[leftmargin=*]
  \item Primary medical condition(s)
  \item Allergies
  \item Current medications
  \item Emergency contact information
  \item Physician/dialysis/contact numbers
\end{itemize}

Always check for identification bracelets/necklaces on unresponsive persons and use this information to guide treatment decisions.

\section{Review Questions}

\begin{enumerate}[leftmargin=*]
\item Why are infants and young children more vulnerable to dehydration and hypothermia?
\item What is the correct positioning for a pregnant woman who needs care, and why?
\item Name three reasons an older adult is at higher risk after a fall.
\item How should you communicate with a person who is hearing impaired during an emergency?
\item Where should you look first for a medical alert ID on an unresponsive person?
\item Why might hand placement differ for CPR on a pregnant person in the later stages?
\end{enumerate}

\index{pediatric first aid}
\index{infant assessment}
\index{child CPR}
\index{geriatric first aid}
\index{stroke in elderly}
\index{fall management}
\index{pregnancy first aid}
\index{left lateral positioning}
\index{visual impairment}
\index{hearing impairment}
\index{mobility impairment}
\index{cognitive disability}
\index{medical alert ID}
